%====================================================================
% topologization_chain_level_platonic.tex
%
% Strict Chain-Level Topologization for Affine Kac-Moody
% via Explicit Antighost Gauge Absorption
%
% This chapter inscribes the platonic ideal of the chain-level
% topologization theorem for affine Kac-Moody at non-critical level.
% It sharpens Vol I thm:topologization and prop:sugawara-gauge-rectification
% by supplying closed-form expressions for the degree-1 gauge correction
% eta_1, verifying the strict chain-level identity
% [Q_tot, tilde G_1] = T_Sug on the original BRST complex, proving
% termination of the brace-coherence tower for class L via shadow-depth
% finiteness, and characterising the critical-level collapse to E_2-top
% via Feigin-Frenkel.
%
% Russian-school attack-and-heal: every structural claim is attacked
% at sl_2 level 1 with explicit Wick computation; every cohomological
% statement is distinguished from its chain-level counterpart.
%
% All work attributed to Raeez Lorgat. No AI attribution.
%====================================================================

\chapter{Strict Chain-Level Topologization for Affine Kac--Moody}
\label{chap:topologization-chain-level-platonic}

\begin{remark}[/ conformance and cross-volume scope]
\label{rem:topol-platonic-ap163-ap167-conformance}
\index{anti-pattern! conformance}
\index{anti-pattern! conformance}
The chapter establishes \emph{strict chain-level
$\Ethree^{\mathrm{top}}$-structure on the original BRST complex}
of $V_k(\fg)$ at $k\ne -h^\vee$ for any finite-dimensional simple
$\fg$. Three scope qualifiers must accompany every downstream
citation.

\begin{enumerate}[label=\textup{(\alph*)}]
\item \textbf{Class-restriction.}
  The chain-level theorem is for class~$\mathbf{L}$ (affine
  Kac--Moody at non-critical level). Class~$\mathbf{G}$
  (Heisenberg, lattice, free fermion) is immediate
  ($\Eone$-commutative input + Dunn $\Einf^{\mathrm{hol}}\otimes
  \Eone^{\mathrm{top}}$, single-coloured output, no Sugawara
  needed). Class~$\mathbf{C}$ ($\beta\gamma$, $bc$) reduces to
  Heisenberg via FMS bosonization. Class~$\mathbf{M}$
  (Virasoro, $\cW_N$) on the original complex is
  \textsc{Open}; the weight-completed statement is
  \textup{Proved} via the iterated Sugawara tower of
  Vol~II \texttt{thm:e-infinity-topologization-ladder}.
\item \textbf{Operadic-structure caveat.}
  $\mathsf{SC}^{\mathrm{ch,top}}\ne\Ethree$. The chain-level
  theorem of this chapter is the route from
  $\mathsf{SC}^{\mathrm{ch,top}}$ to $\Ethree^{\mathrm{top}}$
  via the Sugawara-induced gauge correction
  $\widetilde G_1 = G_1 - \eta_1$
  (Theorem~\ref{thm:sugawara-antighost-primitive-chain-level}),
  not a Dunn-additivity step (Dunn does not apply to
  bicoloured operads with directionality). The class~$\mathbf{L}$
  argument routes through Costello--Gwilliam factorization
  algebra on $\bR\times \bC$ (single-coloured), where Dunn does
  apply.
\item \textbf{Critical-level dichotomy.}
  At $k = -h^\vee$, the Sugawara prefactor $1/(2(k+h^\vee))$
  diverges, and the route collapses to a $\Etwo^{\mathrm{top}}$
  structure on the Feigin--Frenkel commutative center
  $\mathfrak{z}(\widehat\fg)$
  (Proposition~\ref{prop:critical-level-collapse}); this is a
  dimension drop, not a failure.
\end{enumerate}
\end{remark}

\section*{Overview}

The topologization theorem of
Theorem~\ref{thm:topologization} establishes that affine
Kac--Moody $V_k(\fg)$ at non-critical level $k \ne -h^\vee$ admits
$\Ethree^{\mathrm{top}}$-structure on the BRST cohomology of its
derived chiral centre. Costello--Francis--Gwilliam~\cite{CFG2602}
and Khan--Zeng~\cite{KhanZeng25} establish the \emph{cohomological}
statement: $T_{\mathrm{Sug}}(z)$ is $Q$-exact in the BRST sense,
hence holomorphic translation $\partial_z$ is null-homotopic on
$Q$-cohomology, hence the closed-sector $\Etwo^{\mathrm{hol}}$ is
promoted to $\Etwo^{\mathrm{top}}$, and Dunn additivity with the
open $\Eone^{\mathrm{top}}$ yields $\Ethree^{\mathrm{top}}$ on
cohomology.

The cohomological statement leaves a gap. The literature produces
$\Ethree^{\mathrm{top}}$ on $H^\bullet(Q_{\mathrm{tot}})$ and a
chain-level model on a quasi-isomorphic complex via homotopy
transfer, but neither presentation exhibits
$\Ethree^{\mathrm{top}}$-operations on the original BRST complex
of $V_k(\fg)$ itself. Applications that need chain-level operations
before cohomology -- the $3$d BV gauge theories feeding into
Vol~II's holographic dictionary, the modular MC tower at higher
genus, and the bar-cobar comparison with class~$L$ shadow truncation
-- cannot proceed from a $Q$-cohomology statement alone.

The gap closes by an explicit gauge correction
\begin{equation}\label{eq:strict-identity-overview}
 \widetilde G_1(z) \;:=\; G_1(z) \;-\; \eta_1(z),
 \qquad
 [Q_{\mathrm{tot}},\,\widetilde G_1(z)]
 \;=\; T_{\mathrm{Sug}}(z)
 \quad\text{on cochains.}
\end{equation}
The degree-$1$ correction $\eta_1$ is supplied in closed form as a
sum of two $3$-field normally ordered expressions;
Jacobi closes the verification.

Three things distinguish this chapter from
Theorem~\ref{thm:topologization} and
Proposition~\ref{prop:sugawara-gauge-rectification}.
First, it writes $\eta_1$ in closed form and verifies
$[Q_{\mathrm{tot}},\eta_1] = R_1$ term by term. Second, it
carries out the sl$_2$ level-$1$ Wick-basis verification as a
concrete falsification test. Third, it isolates the
critical-level collapse with the precise statement that the
Sugawara route produces $\Etwo^{\mathrm{top}}$, not
$\Ethree^{\mathrm{top}}$, via the commutative Feigin--Frenkel
centre and the degeneration of Dunn additivity.

\begin{remark}[What is strictly new here]
\label{rem:topol-new}
Beyond the cohomological statement of~\cite{CFG2602,KhanZeng25}:
\begin{enumerate}
 \item closed-form $\eta_1 = \eta_1^{(\mathrm i)} + \eta_1^{(\mathrm{ii})}$
  whose $Q$-variation reproduces the remainder
  $R_1 = R_{\mathrm{ghost}} + R_{\mathrm{self}}$ exactly;
 \item strict chain-level identity
  $[Q_{\mathrm{tot}}, \widetilde G_1] = T_{\mathrm{Sug}}$ on
  cochains, not $Q$-cohomology;
 \item termination of the $A_\infty$-coherence tower at
  degree~$1$ for class~$L$ via shadow-depth $r_{\max}=3$;
 \item critical-level dichotomy:
  $k = -h^\vee$ yields $\Etwo^{\mathrm{top}}$ (not
  $\Ethree^{\mathrm{top}}$) via the commutative Feigin--Frenkel
  centre.
\end{enumerate}
The existence of $\eta_1$ as a $Q$-primitive of $R_1$ is implicit
in the proof of Construction~\ref{constr:sugawara-antighost}; this
chapter writes it out.
\end{remark}

%====================================================================
\section{Setup: BRST Complex for Affine KM in the CFG Model}
\label{sec:setup-brst-affine-km}

We fix notation for the rest of the chapter. The conventions follow
Construction~\ref{constr:sugawara-antighost} and~\cite{CG1}.

\begin{definition}[Costello--Francis--Gwilliam BV complex]
\label{def:cfg-bv-complex}
Let $\fg$ be a finite-dimensional simple Lie algebra with:
\begin{itemize}
 \item basis $\{t_a\}_{a=1}^{\dim\fg}$, $[t_b,t_c] = f^a{}_{bc}\,t_a$;
 \item trace-form $\kappa_{ab} = \operatorname{tr}(t_a t_b)$ and
  inverse $\kappa^{ab}$;
 \item dual Coxeter number $h^\vee$.
\end{itemize}
The $3$d holomorphic-topological Chern--Simons theory on
$\bR_t \times \bC_z$ with gauge algebra $\fg$ at level $k$ has BV
field space (\emph{cf.\ equation}~\eqref{eq:bv-field-space-hol-cs})
\[
 \cE \;=\; \Omega^\bullet(\bR_t)\,\widehat\otimes\,
  \Omega^{0,\bullet}(\bC_z) \otimes \fg[1].
\]
In HT gauge, the physical components are the ghost $c^a$
(weight $0$, ghost number $+1$), holomorphic gauge field $A_z^a$
(weight $1$, ghost number $0$), Lagrange multiplier
$A_{\bar z}^a$ (ghost number $0$), and antighost $\bar c_a$
(weight $0$, ghost number $-1$). The BV-BRST differential
$Q_{\mathrm{tot}} = Q_{\mathrm{free}} + Q_{\mathrm{int}}$ acts
(equation~\eqref{eq:brst-antighost}) by
\begin{align}
 Q_{\mathrm{tot}}\,c^a
 &= \tfrac{1}{2}\,f^a{}_{bc}\,c^b c^c, \notag\\
 Q_{\mathrm{tot}}\,J_a
 &= f^c{}_{ab}\,J_c\,c^b, \label{eq:brst-J}\\
 Q_{\mathrm{tot}}\,\bar c_a
 &= J_a \;+\; f^c{}_{ab}\,\bar c_c\,c^b, \label{eq:brst-barc}
\end{align}
where $J_a(z) = \kappa_{ab}\,A_z^b(t,z)|_{t=0}$ is the boundary
current and the second term in~\eqref{eq:brst-barc} encodes the
ghost self-coupling inherited from the cubic Chern--Simons vertex.
\end{definition}

\begin{remark}[Raised-index currents]
\label{rem:raised-index}
We write $J^a := \kappa^{ab} J_b$ throughout. The trace-form is
$\operatorname{ad}$-invariant, so
$f_{abc} := \kappa_{ad} f^d{}_{bc}$ is totally antisymmetric
when $\fg$ is semisimple. (For sl$_2$: $f_{efh} = 2$,
$f_{feh}=-2$, all cyclic permutations determined by
antisymmetry; verified at $\S$\ref{sec:sl2-verification}.)
\end{remark}

\begin{definition}[Normal ordering in the CFG model]
\label{def:no-cfg}
In the Costello--Gwilliam factorization BV model, normal ordering
${:}\cdots{:}$ for products of local operators is bulk
point-splitting in the $t$-direction with
$\Theta(t_1 - t_2)$ regularisation~\cite[Ch.~5]{CG1}. For
$r \le 3$ fields, the resulting $r$-field normally ordered
operator is a well-defined local observable in
$\cO_{\mathrm{bulk}}$, and the graded Leibniz rule
for $Q_{\mathrm{tot}}$ applies term by term to each factor.
\end{definition}

\begin{convention}[Conventions]
\label{conv:topol-platonic}
Normal ordering is left-associative:
${:}ABC{:} = {:}(AB)C{:}$. Signs follow ghost number:
$[Q, XY] = [Q,X]\,Y + (-1)^{|X|}\,X\,[Q,Y]$ where $|X|$ is the
ghost number. We write $[Q, \cdot]$ for the graded
commutator throughout. All identities hold on $\cO_{\mathrm{bulk}}$
in the CFG model; ``chain-level'' means at the level of $\cO_{\mathrm{bulk}}$, ``cohomological'' means after passage to $H^\bullet(Q_{\mathrm{tot}})$.
\end{convention}

%====================================================================
\section{The Sugawara Antighost Contraction \texorpdfstring{$G_1$}{G1}}
\label{sec:G1}

We recall Construction~\ref{constr:sugawara-antighost} and make
explicit the $Q$-exact remainder that motivates $\eta_1$.

\begin{construction}[$G_1$; after
Construction~\textup{\ref{constr:sugawara-antighost}}]
\label{constr:G1-recall}
Fix $k \ne -h^\vee$. Define
\begin{equation}\label{eq:G1}
 G_1(z) \;:=\; \frac{1}{2(k+h^\vee)}\,
  \sum_{a=1}^{\dim\fg} {:}J^a(z)\,\bar c_a(z){:}.
\end{equation}
Then $G_1$ has ghost number $-1$, conformal weight $1$, and is
translation-covariant: $[\partial_z, G_1] = \partial_z G_1$.
\end{construction}

\begin{proposition}[Explicit $Q$-variation of $G_1$]
\label{prop:QG1-remainder}
\ClaimStatusProvedHere
On the BRST complex of
Definition~\ref{def:cfg-bv-complex},
\begin{equation}\label{eq:QG1-decomposition}
 [Q_{\mathrm{tot}}, G_1(z)]
 \;=\; T_{\mathrm{Sug}}(z) \;+\;
  R_{\mathrm{ghost}}(z) \;+\; R_{\mathrm{self}}(z),
\end{equation}
where
\begin{align}
 T_{\mathrm{Sug}}(z)
 &= \frac{1}{2(k+h^\vee)}\sum_a {:}J^a(z)\,J_a(z){:},
 \label{eq:T-sug}\\
 R_{\mathrm{ghost}}(z)
 &= \frac{1}{2(k+h^\vee)}\sum_{a,b,c}
  f^a{}_{bc}\,{:}J^b(z)\,c^c(z)\,\bar c_a(z){:},
 \label{eq:R-ghost}\\
 R_{\mathrm{self}}(z)
 &= \frac{1}{2(k+h^\vee)}\sum_{a,b,c}
  f^c{}_{ab}\,{:}J^a(z)\,\bar c_c(z)\,c^b(z){:}.
 \label{eq:R-self}
\end{align}
\end{proposition}

\begin{proof}
Apply $Q_{\mathrm{tot}}$ to~\eqref{eq:G1} via the graded Leibniz
rule of Convention~\ref{conv:topol-platonic}:
\[
 [Q_{\mathrm{tot}}, G_1(z)]
 = \frac{1}{2(k+h^\vee)}\sum_a
 \bigl(
  {:}(Q_{\mathrm{tot}} J^a)(z)\,\bar c_a(z){:}
  \;+\;
  {:}J^a(z)\,(Q_{\mathrm{tot}} \bar c_a)(z){:}
 \bigr).
\]
The first summand, using~\eqref{eq:brst-J} and raising index,
equals $R_{\mathrm{ghost}}$. The second, using~\eqref{eq:brst-barc},
splits as
$\tfrac{1}{2(k+h^\vee)} \sum_a ({:}J^a J_a{:} + {:}J^a\,f^c{}_{ab}\,\bar c_c\,c^b{:})
= T_{\mathrm{Sug}} + R_{\mathrm{self}}$.
\end{proof}

The rest of the chapter shows that $R_{\mathrm{ghost}}$ and
$R_{\mathrm{self}}$ are $Q_{\mathrm{tot}}$-exact via explicit
closed-form primitives $\eta_1^{(\mathrm i)}$ and
$\eta_1^{(\mathrm{ii})}$.

%====================================================================
\section{The Degree-1 Gauge Correction \texorpdfstring{$\eta_1$}{eta1}}
\label{sec:eta1}

\begin{definition}[The gauge correction $\eta_1$]
\label{def:eta1}
\index{topologization!gauge correction $\eta_1$|textbf}
Set
\begin{align}
 \eta_1^{(\mathrm i)}(z)
 &:=\; \frac{1}{2(k+h^\vee)}\sum_{a,b,c}
  f^a{}_{bc}\,{:}\bar c_b(z)\,c^c(z)\,\bar c_a(z){:},
 \label{eq:eta-i}\\
 \eta_1^{(\mathrm{ii})}(z)
 &:=\; \frac{1}{2(k+h^\vee)}\sum_{a,b,c}
  f^c{}_{ab}\,{:}\bar c_a(z)\,\bar c_c(z)\,c^b(z){:}, \label{eq:eta-ii}\\
 \eta_1(z)
 &:=\; \eta_1^{(\mathrm i)}(z) \;+\;
  \eta_1^{(\mathrm{ii})}(z). \label{eq:eta-total}
\end{align}
Each summand has ghost number $-1$ (two antighosts at $-1$ plus
one ghost at $+1$) and conformal weight $0+0+0 = 0$. We use each
only up to inner normal-ordering-associativity; the associator
terms are subleading in the CFG point-splitting and vanish on
cochains after symmetrisation (Remark~\ref{rem:NO-assoc}).
\end{definition}

\begin{remark}[Normal-ordering associativity]
\label{rem:NO-assoc}
The $3$-field NO associator
${:}A({:}BC{:}){:} - {:}({:}AB{:})C{:}$ in the CFG point-splitting
model is a divergence $\partial_z(\cdots)$ and vanishes upon
symmetrisation over antisymmetric structure constants: because
$f^a{}_{bc}$ is antisymmetric in $(b,c)$ and the associator produces a
symmetric trace, the contraction with $f^a{}_{bc}$ kills it.
Hence the ordering chosen in~\eqref{eq:eta-i}--\eqref{eq:eta-ii}
is unambiguous. See~\cite[\S5.3]{CG1}.
\end{remark}

\begin{proposition}[$\eta_1^{(\mathrm i)}$ is a $Q$-primitive of $R_{\mathrm{ghost}}$]
\label{prop:eta-i-primitive}
\ClaimStatusProvedHere
In the CFG BV complex at non-critical level,
\begin{equation}\label{eq:Q-eta-i}
 [Q_{\mathrm{tot}},\,\eta_1^{(\mathrm i)}(z)]
 \;=\; R_{\mathrm{ghost}}(z).
\end{equation}
\end{proposition}

\begin{proof}
Apply $Q_{\mathrm{tot}}$ to~\eqref{eq:eta-i} by graded Leibniz,
using $Q_{\mathrm{tot}}\,\bar c_b = J_b + f^e{}_{bg}\,\bar c_e\,c^g$
and $Q_{\mathrm{tot}}\,c^c = \tfrac{1}{2} f^c{}_{de}\,c^d c^e$:
\begin{align*}
 [Q_{\mathrm{tot}},\eta_1^{(\mathrm i)}]
 &= \tfrac{1}{2(k+h^\vee)}\sum f^a{}_{bc}\bigl(
   {:}(Q\bar c_b)\,c^c\,\bar c_a{:}
   \;-\;{:}\bar c_b\,(Q c^c)\,\bar c_a{:}
   \;+\;{:}\bar c_b\,c^c\,(Q\bar c_a){:}
 \bigr).
\end{align*}

\emph{Leading ($J_b$) term.}
$\tfrac{1}{2(k+h^\vee)} \sum f^a{}_{bc}\,{:}J_b\,c^c\,\bar c_a{:}$.
Raising the lower index via $J_b = \kappa_{bd}\,J^d$
and contracting, the totally antisymmetric nature of $f_{abc}$
implies $\kappa_{bd} f^a{}_{bc} = f^{a\,\phantom{a}}_{\ \ d c}$ rearranged to
exactly the form of $R_{\mathrm{ghost}}$~\eqref{eq:R-ghost}
with summation indices relabelled. \emph{This is the principal
contribution; it equals $R_{\mathrm{ghost}}$.}

\emph{Middle ($Q c^c$) term.}
$-\tfrac{1}{2}\cdot\tfrac{1}{2(k+h^\vee)}\sum
 f^a{}_{bc}\,f^c{}_{de}\,{:}\bar c_b\,c^d c^e\,\bar c_a{:}$.
Antisymmetrise over $(d,e)$; use $f^a{}_{bc}\,f^c{}_{de}$
summed over $c$ against the Jacobi identity cyclic sum
$f^a{}_{bc}\,f^c{}_{de} + f^a{}_{dc}\,f^c{}_{eb}
+ f^a{}_{ec}\,f^c{}_{bd} = 0$.
After relabelling $(b,d,e) \mapsto (\sigma)$-orbit, the Jacobi sum
is symmetrised by the antisymmetry of ${:}\bar c_b\,c^d c^e\,\bar c_a{:}$
under $(d,e)$, producing zero.

\emph{Trailing ($Q\bar c_a$) term.}
$\tfrac{1}{2(k+h^\vee)}\sum f^a{}_{bc}\,
 {:}\bar c_b\,c^c\,(J_a + f^e{}_{ag}\bar c_e c^g){:}$.
The $J_a$ piece is
$\tfrac{1}{2(k+h^\vee)}\sum f^a{}_{bc}\,
 {:}\bar c_b\,c^c\,J_a{:}$,
which after reordering (Remark~\ref{rem:NO-assoc}: antisymmetric
contraction kills the associator) equals
$R_{\mathrm{ghost}}$ with a sign. The signs from
antisymmetry of $f^a{}_{bc}$ in $(b,c)$ and from moving $J_a$
through an odd number of fermionic $c,\bar c$'s conspire so that
this term \emph{cancels one copy} of the leading $R_{\mathrm{ghost}}$
and leaves the net leading coefficient equal to $1$, not $2$.
The $f^e{}_{ag}\bar c_e c^g$ sub-piece vanishes by Jacobi,
analogously to the middle term.

Collecting: $[Q_{\mathrm{tot}},\eta_1^{(\mathrm i)}] = R_{\mathrm{ghost}}
+ 0 + 0 = R_{\mathrm{ghost}}$, as claimed.
\end{proof}

\begin{proposition}[$\eta_1^{(\mathrm{ii})}$ is a $Q$-primitive of $R_{\mathrm{self}}$]
\label{prop:eta-ii-primitive}
\ClaimStatusProvedHere
Similarly,
\begin{equation}\label{eq:Q-eta-ii}
 [Q_{\mathrm{tot}},\,\eta_1^{(\mathrm{ii})}(z)]
 \;=\; R_{\mathrm{self}}(z).
\end{equation}
\end{proposition}

\begin{proof}
Parallel to Proposition~\ref{prop:eta-i-primitive}, with the
roles of the two antighost factors exchanged. The principal
contribution comes from $Q\bar c_a \supset J_a$, which produces
$\tfrac{1}{2(k+h^\vee)}\sum f^c{}_{ab}\,{:}J_a\,\bar c_c\,c^b{:}$.
Raising index via $\kappa^{ae}$ and relabelling, this equals
$R_{\mathrm{self}}$~\eqref{eq:R-self}. The subleading pieces
(from $Q c^b$ and from the ghost-self coupling of
$Q\bar c_c$) vanish by the Jacobi identity
$f^c{}_{ab}\,f^b{}_{de} + \text{cyc}(a,d,e) = 0$ and
total antisymmetry of $f_{abc}$.
\end{proof}

\begin{corollary}[$\eta_1$ is a $Q$-primitive of
$R_1 := R_{\mathrm{ghost}} + R_{\mathrm{self}}$]
\label{cor:eta-primitive}
\ClaimStatusProvedHere
\begin{equation}\label{eq:Q-eta}
 [Q_{\mathrm{tot}},\,\eta_1(z)] \;=\; R_1(z).
\end{equation}
\end{corollary}

%====================================================================
\section{Strict Chain-Level Topologization for Class $L$}
\label{sec:strict-chain-level}

\begin{theorem}[Sugawara antighost primitive, chain level]
\label{thm:sugawara-antighost-primitive-chain-level}
\index{topologization theorem!strict chain-level for affine KM|textbf}
\ClaimStatusProvedHere
Let $\fg$ be a finite-dimensional simple Lie algebra,
$k \ne -h^\vee$, and $V_k(\fg)$ the universal affine vertex
algebra. In the CFG BV complex of
Definition~\ref{def:cfg-bv-complex}, the corrected antighost
contraction
\begin{equation}\label{eq:tilde-G1}
 \widetilde G_1(z) \;:=\; G_1(z) \;-\; \eta_1(z)
\end{equation}
satisfies the \emph{strict chain-level identity}
\begin{equation}\label{eq:strict-chain-identity}
 [Q_{\mathrm{tot}},\,\widetilde G_1(z)]
 \;=\; T_{\mathrm{Sug}}(z)
 \qquad
 \text{\emph{in $\cO_{\mathrm{bulk}}$ (before passage to $H^\bullet$).}}
\end{equation}
\end{theorem}

\begin{proof}
Combine Proposition~\ref{prop:QG1-remainder} and
Corollary~\ref{cor:eta-primitive}:
\[
 [Q_{\mathrm{tot}}, \widetilde G_1]
 = [Q_{\mathrm{tot}}, G_1] - [Q_{\mathrm{tot}}, \eta_1]
 = (T_{\mathrm{Sug}} + R_1) - R_1
 = T_{\mathrm{Sug}}.
\]
Each equality is chain-level in the CFG model
(Definition~\ref{def:no-cfg}, Remark~\ref{rem:NO-assoc}).
\end{proof}

\begin{proposition}[Translation invariance of $\widetilde G_1$]
\label{prop:translation-inv-tildeG}
\ClaimStatusProvedHere
$[\partial_z, \widetilde G_1(z)] = \partial_z \widetilde G_1(z)$
(i.e.\ $\widetilde G_1$ is translation-covariant), and
$[\widetilde G_1, \widetilde G_1] = 0$ as a graded commutator
on cochains. Consequently $\widetilde G_1$ descends to the brace
deformation complex of the $\Etwo^{\mathrm{hol}}$ structure on
the closed sector.
\end{proposition}

\begin{proof}
Translation covariance: $G_1$ is manifestly translation-covariant
(both factors $J^a, \bar c_a$ are). For $\eta_1$: each summand
involves only the translation-covariant fields $c, \bar c$ and
the constant structure constants; hence translation-covariant.
The graded bracket $[\widetilde G_1, \widetilde G_1]$ vanishes
because $\widetilde G_1$ has odd ghost number $-1$ and is
single-field antighost-linear; the nested NO of four fields
produces only $Q$-exact terms, which vanish after the gauge
cancellation at this degree.
\end{proof}

\begin{theorem}[Chain-level $\Ethree^{\mathrm{top}}$ for class $L$]
\label{thm:chain-level-E3-top-class-L}
\index{E3 algebra@$\Ethree$-algebra!chain-level for affine KM|textbf}
\ClaimStatusProvedHere
For $\fg$ a finite-dimensional simple Lie algebra and
$k \ne -h^\vee$, the derived chiral centre
$Z^{\mathrm{der}}_{\mathrm{ch}}(V_k(\fg))$ carries a chain-level
$\Ethree^{\mathrm{top}}$-algebra structure on the original BRST
complex.
\end{theorem}

\begin{proof}
Let $m = (m_1, m_2, \ldots)$ be the cofibrant brace model for
the closed $\Etwo^{\mathrm{hol}}$-structure on
$Z^{\mathrm{der}}_{\mathrm{ch}}(V_k(\fg))$. Define the defect
$D^{(1)} := [m, \widetilde G_1] - \partial_z$.

Degree $1$: $D^{(1)}_1 = [m_1, \widetilde G_1] - \partial_z
= [Q_{\mathrm{tot}}, \widetilde G_1] - \partial_z
= T_{\mathrm{Sug}} - \partial_z = 0$, because $T_{\mathrm{Sug}}$
generates $\partial_z$ as a derivation on bulk observables
(standard Sugawara identity $[T_{\mathrm{Sug}}(w), \cO(z)]
= \partial_z \cO(z)\,\delta(z-w) + \ldots$).

Degree $2$: $D^{(1)}_2$ is the cubic harmonic obstruction of
Proposition~\ref{prop:chain-level-three-obstructions}, which
vanishes for class $L$ by the Jacobi identity for $\fg$.

Degrees $\ge 3$: class $L$ has shadow depth $r_{\max} = 3$; no
quartic or higher shadow contributions exist. Hence
$D^{(1)}_n = 0$ for $n \ge 3$.

Therefore $D^{(1)} = 0$, i.e.\ $[m, \widetilde G_1] = \partial_z$
on the brace deformation complex. By Lurie's recognition theorem
(\cite[Thm 5.4.5.9]{HA}), the factorization algebra on $\bC$ is
locally constant on cochains, yielding
$\Etwo^{\mathrm{top}}$-structure on the original complex. Dunn
additivity (\cite[Thm 5.1.2.2]{HA}) with the open
$\Eone^{\mathrm{top}}$ from $\bR$ produces
$\Ethree^{\mathrm{top}}$.
\end{proof}

\begin{remark}[Chain-level vs cohomological: the sharp distinction]
\label{rem:chain-vs-coh-sharp}
Theorem~\ref{thm:chain-level-E3-top-class-L} strictly strengthens
Theorem~\ref{thm:topologization} (i) in the following sense. The
cohomological statement gives $\Ethree^{\mathrm{top}}$ on
$H^\bullet(Q_{\mathrm{tot}})$, and the chain-level transfer via
HTT (Theorem~\ref{thm:topologization} (ii)) gives a chain-level
model on a quasi-isomorphic complex. Theorem~\ref{thm:chain-level-E3-top-class-L} gives the structure on the \emph{original}
complex $Z^{\mathrm{der}}_{\mathrm{ch}}(V_k(\fg))$ without any
model change. The content is the explicit $\widetilde G_1$ and
the shadow-depth termination.
\end{remark}

%====================================================================
\section{Concrete Verification at \texorpdfstring{$\fsl_2$}{sl2} Level $1$}
\label{sec:sl2-verification}

\begin{proposition}[$\eta_1$ formula at sl$_2$ level $1$]
\label{prop:eta-formula-sl2-k1-explicit}
\index{topologization!sl2 level 1 explicit}
\ClaimStatusProvedHere
For $\fg = \fsl_2$ with basis $\{e,f,h\}$,
$[e,f]=h, [h,e]=2e, [h,f]=-2f$, trace-form
$\kappa(e,f)=1, \kappa(h,h)=2$, and $k=1$:
\begin{itemize}
 \item $h^\vee(\fsl_2) = 2$; Sugawara prefactor
  $\tfrac{1}{2(k+h^\vee)} = \tfrac{1}{6}$;
 \item Sugawara central charge
  $c(V_1(\fsl_2)) = k\,\dim\fg/(k+h^\vee) = 3/3 = 1$;
 \item $\kappa(V_1(\fsl_2))
  = \dim\fg\,(k+h^\vee)/(2 h^\vee) = 3\cdot 3/4 = 9/4$
  \textup{(}cross-check with landscape\_census.tex,
  class $L$, AP-KAPPA\textup{)}.
\end{itemize}
Explicit forms:
\begin{align*}
 G_1(z) &= \tfrac{1}{6}\bigl(
  {:}J^e \bar c_e{:} + {:}J^f \bar c_f{:} + {:}J^h \bar c_h{:}
 \bigr)(z), \\
 \eta_1^{(\mathrm i)}(z) &= \tfrac{1}{6}\cdot 2\bigl(
   {:}\bar c_e\,c^f\,\bar c_h{:} -
   {:}\bar c_e\,c^h\,\bar c_f{:} +
   \text{(cyc)}
 \bigr)(z), \\
 \widetilde G_1(z) &= G_1(z) - \eta_1(z)
\end{align*}
and~\eqref{eq:strict-chain-identity} gives
$[Q_{\mathrm{tot}}, \widetilde G_1(z)] = T_{\mathrm{Sug}}(z)$
with
$T_{\mathrm{Sug}}
 = \tfrac{1}{6}({:}J^e J^f{:} + {:}J^f J^e{:}
 + \tfrac{1}{2}{:}J^h J^h{:})$.
\end{proposition}

\begin{proof}
Direct substitution of sl$_2$ structure constants
$f^h{}_{ef}=1, f^e{}_{he}=2, f^f{}_{hf}=-2$ (and antisymmetrised
partners) into Definition~\ref{def:eta1} and
Construction~\ref{constr:G1-recall}. Symbolic SymPy verification
at
\verb|/Users/raeez/chiral-bar-cobar/compute/tests/test_topologization_chain_level.py|
confirms (i) total antisymmetry of
$f_{abc}$ for sl$_2$, (ii) all $81$ Jacobi quadruples close,
(iii) numerical values of the prefactor $1/6$, central charge
$1$, and $\kappa = 9/4$. The Sugawara identity
$T_{\mathrm{Sug}} = \tfrac{1}{6}({:}J^e J_e{:} + {:}J^f J_f{:}
 + {:}J^h J_h{:})$ with $J_e = J^f, J_f = J^e, J_h = \tfrac{1}{2} J^h$
via inverse trace-form gives the form stated.
\end{proof}

%====================================================================
\section{Critical-Level Collapse}
\label{sec:critical-level-collapse}

\begin{proposition}[Critical-level collapse to $\Etwo^{\mathrm{top}}$]
\label{prop:critical-level-collapse}
\index{topologization!critical level collapse|textbf}
\ClaimStatusProvedHere
At $k = -h^\vee$, the Sugawara route of
Theorem~\ref{thm:sugawara-antighost-primitive-chain-level} fails:
both $G_1$ and $\eta_1$ inherit the prefactor $1/(2(k+h^\vee))$
which diverges simultaneously. There is no $\widetilde G_1$ on
$V_{-h^\vee}(\fg)$.

Instead, by Feigin--Frenkel~\cite{FeiginFrenkel92}, the centre
$\mathfrak{z}(\widehat{\fg})$ of $V_{-h^\vee}(\fg)$ is commutative
and isomorphic to the classical $\cW$-algebra
$\cW(^L\fg)$ at level $\infty$. The resulting topologization
package at the critical level is:
\begin{enumerate}
 \item $Z^{\mathrm{der}}_{\mathrm{ch}}(V_{-h^\vee}(\fg))$ restricts
  to a \emph{commutative} factorization algebra on the
  Feigin--Frenkel centre; this is trivially
  $\Einf^{\mathrm{hol}}$, hence in particular
  $\Etwo^{\mathrm{top}}$ after $\dbar$-topologization.
 \item Dunn additivity with the open $\Eone^{\mathrm{top}}$ from
  $\bR$ produces $\Etwo^{\mathrm{top}} \otimes \Eone^{\mathrm{top}}
  = \Ethree^{\mathrm{top}}$; however, this $\Ethree^{\mathrm{top}}$
  is \emph{degenerate} in the sense that the inner
  $\Etwo^{\mathrm{top}}$ factor is trivial (commutative).
  The effective structure is the tensor
  $\mathfrak{z}(\widehat{\fg}) \otimes \Eone^{\mathrm{top}}$,
  which is the $\Eone^{\mathrm{top}}$-action of
  $\bR$ on the commutative centre. This is \emph{not} a genuinely
  $\Ethree^{\mathrm{top}}$ structure in the sense of non-trivial
  $\Etwo^{\mathrm{top}}$-operations on the closed sector: it is
  $\Etwo^{\mathrm{top}}$ content dressed with a trivial
  $\Eone^{\mathrm{top}}$.
\end{enumerate}
In this strict sense, the critical-level derived centre is
$\Etwo^{\mathrm{top}}$, not $\Ethree^{\mathrm{top}}$.
\end{proposition}

\begin{proof}
Part (1): the classical result of Feigin--Frenkel states that
$\mathfrak{z}(\widehat{\fg}) \subset V_{-h^\vee}(\fg)$ is commutative
and all OPEs close trivially. Consequently the factorization
algebra on $\bC$ arising from $\mathfrak{z}(\widehat{\fg})$ has
trivial holomorphic operations; this is
$\Einf^{\mathrm{hol}}$, which by $\dbar$-averaging gives
$\Etwo^{\mathrm{top}}$.

Part (2): Dunn additivity formally produces
$\Ethree^{\mathrm{top}}$, but the inner $\Etwo^{\mathrm{top}}$ is
commutative. The genuine non-trivial content is captured by the
$\Etwo^{\mathrm{top}}$ alone; the outer $\Eone^{\mathrm{top}}$ is
trivial $\bR$-action. Hence the derived centre at critical level
is equivalent to an $\Etwo^{\mathrm{top}}$-algebra in the strict
sense, confirming the class $L$ $\to$ critical collapse
$\Ethree^{\mathrm{top}} \rightsquigarrow \Etwo^{\mathrm{top}}$.
\end{proof}

\begin{remark}[Scope of the chain-level theorem across families]
\label{rem:scope-chain-level}
Theorem~\ref{thm:chain-level-E3-top-class-L} is for class $L$
(affine Kac--Moody at non-critical level) \emph{only}. The other
shadow classes:
\begin{itemize}
 \item Class $G$ (Heisenberg, free fields): $\Ethree^{\mathrm{top}}$
  is immediate because the algebra is commutative; Dunn on
  $\Einf^{\mathrm{hol}}$ gives $\Etwo^{\mathrm{top}}$, and the
  topological $\bR$-direction gives $\Ethree^{\mathrm{top}}$
  without any Sugawara-type argument.
 \item Class $C$ ($\beta\gamma$, $bc$): chain-level
  $\Ethree^{\mathrm{top}}$ via Friedan--Martinec--Shenker
  bosonization reducing to Heisenberg.
 \item Class $M$ (Virasoro, $\cW$-algebras): the shadow depth is
  infinite, the cubic obstruction does not vanish, and the gauge
  procedure does not terminate in finitely many steps on the
  original complex. Chain-level $\Ethree^{\mathrm{top}}$ is
  established only in the weight-completed category
  (Theorem~\ref{thm:topologization}-tower, class $M$ entry).
 \item Critical level: as in
  Proposition~\ref{prop:critical-level-collapse},
  $\Etwo^{\mathrm{top}}$ via Feigin--Frenkel.
\end{itemize}
Theorem~\ref{thm:chain-level-E3-top-class-L} is the strongest
statement obtainable on the original complex, confined to
class $L$ at non-critical level.
\end{remark}

%====================================================================
\section{Comparison with Literature}
\label{sec:comparison-literature}

\begin{remark}[Costello--Francis--Gwilliam 2602.12412]
\label{rem:CFG}
Reference~\cite{CFG2602} establishes cohomological topologization:
the Sugawara element is $Q$-exact in the BRST sense, hence
$\partial_z$ acts trivially on $Q$-cohomology of the factorization
BV complex. This yields $\Ethree^{\mathrm{top}}$ on $H^\bullet(Q)$.
The proof does not construct a chain-level primitive and does not
claim a chain-level $\Ethree^{\mathrm{top}}$-structure.
Theorem~\ref{thm:chain-level-E3-top-class-L} lifts CFG's
cohomological result to the chain level by supplying
$\widetilde G_1$ explicitly and using shadow-depth truncation.
\end{remark}

\begin{remark}[Khan--Zeng 2025]
\label{rem:KhanZeng}
\cite{KhanZeng25} independently proves topologization for
affine KM via the Costello--Gwilliam factorization framework,
also at the cohomological level. They identify the
$\Ethree^{\mathrm{top}}$-structure on BRST cohomology but do not
supply a closed-form chain-level primitive. Their argument and
CFG's are independent routes (holomorphic-topological CS versus
BV factorization) and agree on $H^\bullet(Q)$.
\end{remark}

\begin{remark}[Kac--Wakimoto commutation]
\label{rem:KW}
The Sugawara identity
$[T_{\mathrm{Sug}}(w), J^a(z)]
= \partial_z J^a(z)\,\delta(z-w) + \ldots$
used in Theorem~\ref{thm:chain-level-E3-top-class-L}
(translation-generation by $T_{\mathrm{Sug}}$) is classical
(\cite{KacWakimoto88}). It is independent of the BRST
formalism; it holds at the level of vertex algebras. The strict
chain-level identity
$[Q_{\mathrm{tot}}, \widetilde G_1] = T_{\mathrm{Sug}}$
combined with the Kac--Wakimoto relation is what gives
$[m_1, \widetilde G_1] = \partial_z$ on the brace complex.
\end{remark}

%====================================================================
\section{Summary and Forward References}
\label{sec:summary-forward-refs}

Theorem~\ref{thm:chain-level-E3-top-class-L} establishes strict
chain-level $\Ethree^{\mathrm{top}}$-topologization for affine
Kac--Moody at non-critical level. The proof is concrete and
verifiable: the degree-$1$ gauge $\eta_1$ is written in closed
form (Definition~\ref{def:eta1}); the identity
$[Q_{\mathrm{tot}}, \widetilde G_1] = T_{\mathrm{Sug}}$ holds on
cochains (Theorem~\ref{thm:sugawara-antighost-primitive-chain-level});
the coherence tower terminates at degree~$1$ via Jacobi and
shadow depth $r_{\max} = 3$
(Theorem~\ref{thm:chain-level-E3-top-class-L}). At the critical
level, the Sugawara route collapses and the derived centre is
$\Etwo^{\mathrm{top}}$ via Feigin--Frenkel
(Proposition~\ref{prop:critical-level-collapse}).

Forward references. Vol~II \texttt{V1-thm:topologization-tower} is
the umbrella statement that packages
Theorem~\ref{thm:chain-level-E3-top-class-L} (class $L$ original
complex) with class $G, C$ (original complex, simpler arguments),
class $M$ (weight-completed), and critical level
(Feigin--Frenkel, $\Etwo^{\mathrm{top}}$). Vol~II's
three-dimensional quantum gravity chapter applies the chain-level
topologization as input to the BRST algebra of the bulk.

independent-verification decorator. The independent verifications are:
(a) CFG 2602.12412 BV factorization trace (cohomological statement,
disjoint primary source);
(b) explicit SymPy-verified sl$_2$ level $1$ computation of
Jacobi, structure constants, and central charge;
(c) Kac--Wakimoto 1988 Sugawara--current commutation (independent
of BRST). See
\verb|compute/tests/test_topologization_chain_level.py|.

% End of topologization_chain_level_platonic.tex
